\section{Installation du simulateur}

\noindent Cette section décrit les différentes étapes nécessaires pour installer le simulateur sur une machine locale. Les procédures présentées ont été testées sous macOS et Windows.

\subsection{Prérequis}

\noindent Avant d'installer le simulateur, les logiciels suivants doivent être installés sur la machine :

\begin{itemize}
    \item Python (version 3.11 ou supérieure) ;
    \item Git, utilisé pour cloner le dépôt du projet ;
    \item un environnement virtuel Python.
\end{itemize}

\subsection{Installation}

\noindent Le projet est distribué sous la forme d'un dépôt Git. La première étape consiste donc à se placer dans le dossier souhaité, puis à cloner le dépôt sur la machine locale :

\begin{codebox}{bash}
    git clone https://github.com/sayansuos/baggage-handling-rl.git
    cd baggage-handling-rl
\end{codebox}
\vspace{10pt}


\noindent Il est ensuite recommandé de créer un environnement virtuel afin d'isoler les dépendances du projet. Celui-ci peut ensuite être activé selon le système d'exploitation utilisé :

\vspace{10pt}


\begin{itemize}
    \item Sous macOS ou Linux :
          \begin{codebox}{bash}
              python -m venv .venv
              source ./.venv/bin/activate
          \end{codebox}
    \item Sous Windows :
          \begin{codebox}{bash}
              python -m venv .venv
              .\.venv\Scripts\activate
          \end{codebox}
\end{itemize}

\vspace{10pt}

\noindent Une fois l'environnement virtuel activé, les dépendances du projet peuvent être installées à partir du fichier \texttt{requirements.txt} :

\begin{codebox}{bash}
    python -m pip install -r requirements.txt
\end{codebox}
\vspace{10pt}

\noindent Une fois cette étape terminée, le simulateur est prêt à être utilisé.
